% =====================================================================
%  Προοίμιο: γλώσσες, γραμματοσειρές, σελίδα, εικόνες, κώδικας, βιβλιογραφία
% =====================================================================

% ---------------- Λειτουργία final -----------------------------------
% \FinalVersiontrue κρύβει τις σημειώσεις καθοδήγησης. Άλλαξέ το στο τέλος.
\newif\ifFinalVersion
\FinalVersionfalse

% ---------------- Γλώσσες και γραμματοσειρές (XeLaTeX) ----------------
\usepackage{fontspec}
\usepackage{polyglossia}
\setmainlanguage{greek}
\setotherlanguage{english}
% Libertinus: πλήρη ελληνικά, μαθηματικά που ταιριάζουν, διαθέσιμη σε TeX Live και Overleaf.
\setmainfont{Libertinus Serif}
\setsansfont{Libertinus Sans}
\setmonofont{DejaVu Sans Mono}[Scale=MatchLowercase]
\newfontfamily\greekfont{Libertinus Serif}
\newfontfamily\greekfontsf{Libertinus Sans}
\newfontfamily\greekfonttt{DejaVu Sans Mono}[Scale=MatchLowercase]

% ---------------- Μαθηματικά -----------------------------------------
\usepackage{amsmath,amssymb}
\usepackage{unicode-math}
\setmathfont{Libertinus Math}
\usepackage{siunitx}
\sisetup{output-decimal-marker={,}, group-separator={.}, detect-all}

% ---------------- Σελίδα ---------------------------------------------
\usepackage[a4paper,top=2.5cm,bottom=2.5cm,left=3cm,right=2.5cm]{geometry}
\usepackage{setspace}
\onehalfspacing
\usepackage{fancyhdr}
\pagestyle{fancy}
\fancyhf{}
\fancyhead[L]{\small\nouppercase{\leftmark}}
\fancyfoot[C]{\thepage}
\renewcommand{\headrulewidth}{0.4pt}
\fancypagestyle{plain}{\fancyhf{}\fancyfoot[C]{\thepage}\renewcommand{\headrulewidth}{0pt}}
\usepackage{titlesec}
\titleformat{\chapter}[display]{\normalfont\huge\bfseries}{\chaptertitlename\ \thechapter}{16pt}{\Huge}
\titlespacing*{\chapter}{0pt}{0pt}{30pt}

% ---------------- Εικόνες, πίνακες, διαγράμματα ----------------------
\usepackage{graphicx}
\graphicspath{{../figures/}{../figures/generated/}}
\usepackage{float}
\usepackage{booktabs,tabularx,multirow,longtable}
\usepackage{adjustbox}
\usepackage[font=small,labelfont=bf]{caption}
\usepackage{subcaption}
\usepackage{tikz}
\usetikzlibrary{arrows.meta,positioning,fit,backgrounds,calc,shapes.geometric}
\usepackage{pgfplots}
\pgfplotsset{compat=1.18}
\usepackage{xcolor}

% ---------------- Κώδικας --------------------------------------------
\usepackage{listings}
\lstdefinestyle{code}{
  basicstyle=\ttfamily\footnotesize,
  keywordstyle=\color{blue!70!black}\bfseries,
  commentstyle=\color{green!40!black}\itshape,
  stringstyle=\color{red!60!black},
  numbers=left, numberstyle=\tiny\color{gray}, numbersep=6pt,
  frame=single, rulecolor=\color{gray!40},
  breaklines=true, showstringspaces=false, tabsize=4,
  captionpos=b, xleftmargin=1.5em, framexleftmargin=1.5em
}
\lstset{style=code}
\renewcommand{\lstlistingname}{Απόσπασμα}
\renewcommand{\lstlistlistingname}{Κατάλογος αποσπασμάτων κώδικα}

% ---------------- Αλγόριθμοι -----------------------------------------
\usepackage{algorithm}
\usepackage{algpseudocode}
\floatname{algorithm}{Αλγόριθμος}
\renewcommand{\listalgorithmname}{Κατάλογος αλγορίθμων}

% ---------------- Βιβλιογραφία ---------------------------------------
\usepackage{csquotes}
\usepackage[backend=biber,style=numeric-comp,sorting=none,maxbibnames=6,giveninits=true]{biblatex}
\addbibresource{../bibliography/references.bib}

% ---------------- Υπερσύνδεσμοι (τελευταίο) --------------------------
\usepackage[hidelinks,unicode]{hyperref}
\usepackage[nameinlink,noabbrev]{cleveref}
\crefname{chapter}{κεφάλαιο}{κεφάλαια}
\Crefname{chapter}{Κεφάλαιο}{Κεφάλαια}
\crefname{section}{ενότητα}{ενότητες}
\Crefname{section}{Ενότητα}{Ενότητες}
\crefname{subsection}{ενότητα}{ενότητες}
\Crefname{subsection}{Ενότητα}{Ενότητες}
\crefname{figure}{σχήμα}{σχήματα}
\Crefname{figure}{Σχήμα}{Σχήματα}
\crefname{table}{πίνακας}{πίνακες}
\Crefname{table}{Πίνακας}{Πίνακες}
\crefname{equation}{εξίσωση}{εξισώσεις}
\Crefname{equation}{Εξίσωση}{Εξισώσεις}
\crefname{lstlisting}{απόσπασμα}{αποσπάσματα}
\Crefname{lstlisting}{Απόσπασμα}{Αποσπάσματα}
\crefname{algorithm}{αλγόριθμος}{αλγόριθμοι}
\Crefname{algorithm}{Αλγόριθμος}{Αλγόριθμοι}

% ---------------- Σημειώσεις καθοδήγησης -----------------------------
% \guide{...}   : τι γράφεται σε αυτή την ενότητα (μπλε πλαίσιο)
% \todo{...}    : κάτι που λείπει (πορτοκαλί πλαίσιο)
% \fillin{...}  : στοιχείο που πρέπει να συμπληρωθεί (κόκκινο κείμενο)
% \code{...}    : όνομα αρχείου / κλάσης του project
\usepackage{tcolorbox}
\ifFinalVersion
  \newcommand{\guide}[1]{}
  \newcommand{\todo}[1]{}
  \newcommand{\fillin}[1]{#1}
\else
  \newtcolorbox{guidebox}{colback=blue!4,colframe=blue!45!black,boxrule=0.4pt,
    left=4pt,right=4pt,top=2pt,bottom=2pt,fontupper=\small\sffamily,
    title={\scriptsize\sffamily ΟΔΗΓΙΑ},fonttitle=\bfseries,coltitle=white}
  \newtcolorbox{todobox}{colback=orange!8,colframe=orange!70!black,boxrule=0.4pt,
    left=4pt,right=4pt,top=2pt,bottom=2pt,fontupper=\small\sffamily,
    title={\scriptsize\sffamily ΕΚΚΡΕΜΕΙ},fonttitle=\bfseries,coltitle=white}
  \newcommand{\guide}[1]{\begin{guidebox}#1\end{guidebox}}
  \newcommand{\todo}[1]{\begin{todobox}#1\end{todobox}}
  \newcommand{\fillin}[1]{\textcolor{red}{[#1]}}
\fi
\newcommand{\code}[1]{\texttt{\detokenize{#1}}}
\newcommand{\en}[1]{\textenglish{#1}}
